% ============================================================================
% General Conclusion
% ============================================================================
\chapter*{General Conclusion}
\addcontentsline{toc}{chapter}{General Conclusion}
\markboth{General Conclusion}{}

This end-of-year project resulted in the design, development, and deployment
of \leaklens{}, an OSINT platform for breach-data research. The work covered
requirements analysis, security-oriented architecture, and production
deployment on Vercel and a VPS.

The platform enables analysts to query compromised records across seven search
types, enriched with public intelligence APIs. Its security model combines
JWT authentication, RBAC, invite-only registration, rate limiting, and HTTP
hardening (CSP, HSTS, CORS)---applied as defense-in-depth rather than isolated
controls.

Beyond the technical deliverable, this project consolidated skills directly
relevant to cybersecurity practice:

\begin{itemize}[nosep]
  \item Full-stack development with Next.js and FastAPI;
  \item Application security: JWT, RBAC, anti-enumeration, rate limiting;
  \item Containerization with Alpine Docker builds and Trivy scanning;
  \item External API orchestration (Snusbase, LeakCheck, EmailRep, CertSpotter);
  \item Split deployment across Vercel and a VPS with Nginx Proxy Manager.
\end{itemize}

\section*{Future Improvements}

\begin{enumerate}[nosep]
  \item Custom HTTPS domain with Let's Encrypt via Nginx Proxy Manager;
  \item Local breach-data indexing with Elasticsearch;
  \item Automated CI/CD pipeline (GitHub Actions);
  \item TOTP-based two-factor authentication for analyst accounts;
  \item SIEM integration for search audit logging;
  \item Tor-based monitoring for paste sites and darknet forums.
\end{enumerate}

\vspace{0.8cm}
\begin{center}
\parbox{0.82\textwidth}{\itshape\centering
This project shows that combining OSINT capability with rigorous application
security is achievable at student scale---and that understanding both sides
of the problem is essential for any cybersecurity professional.}
\end{center}
